\documentclass[a4paper, 10pt]{article}
\usepackage[utf8]{inputenc}
\usepackage[T1]{fontenc}
\usepackage[spanish]{babel}
\usepackage[top=0.7cm, bottom=0.7cm, left=1.5cm, right=1.5cm]{geometry}
\usepackage{xcolor}
\usepackage{graphicx}
\usepackage[most]{tcolorbox}
\usepackage{parskip}
\usepackage{microtype}
\usepackage{lmodern}
\usepackage{enumitem}
\usepackage{tikz}
\usetikzlibrary{arrows.meta, positioning}

\renewcommand{\familydefault}{\sfdefault}
\pagestyle{empty}
\setlength{\parindent}{0pt}
\setlength{\parskip}{0pt}

%% Paleta de colores YariNET
\definecolor{yblue}{HTML}{1B2A6B}
\definecolor{yteal}{HTML}{00B09B}
\definecolor{ygold}{HTML}{F5C518}
\definecolor{ylight}{HTML}{EEF2FF}
\definecolor{ypurple}{HTML}{7C3AED}
\definecolor{yorange}{HTML}{EA580C}
\definecolor{ygray}{HTML}{475569}
\definecolor{ylightbg}{HTML}{F8FAFC}

\begin{document}

%% ─── HEADER ───────────────────────────────────────────────────────────
\begin{tcolorbox}[
  enhanced,
  colback=yblue,
  colframe=yblue,
  arc=4mm,
  boxrule=0pt,
  left=10pt, right=10pt, top=7pt, bottom=7pt
]
\begin{minipage}[c]{0.20\textwidth}
  \includegraphics[width=\linewidth]{C:/Trabajo/CUMBRE_PLATFORM/Logotipos/YariNet.png}
\end{minipage}%
\hspace{14pt}%
\begin{minipage}[c]{0.76\textwidth}
  {\color{white}\fontsize{30}{36}\selectfont\textbf{YariNET}}\\[3pt]
  {\color{ygold}\fontsize{12.5}{15}\selectfont\textbf{Piensa. Conecta. Act\'{u}a.}}\\[6pt]
  {\color{white!75!yblue}\fontsize{9.5}{12}\selectfont
    Plataforma de deliberaci\'{o}n c\'{i}vica con inteligencia artificial para escolares}\\[2pt]
  {\color{white!50!yblue}\fontsize{8.5}{11}\selectfont
    Educat\'{o}n Challenge 2026\ $\cdot$\ PUCP\ $\cdot$\ Barrio PUCP}
\end{minipage}
\end{tcolorbox}

\vspace{4pt}

%% ─── DESCRIPCIÓN ──────────────────────────────────────────────────────
\begin{tcolorbox}[
  enhanced,
  colback=ylight,
  colframe=yteal,
  arc=3mm,
  boxrule=1.5pt,
  left=12pt, right=12pt, top=7pt, bottom=7pt
]
{\color{yblue}\fontsize{10.5}{14}\selectfont
YariNET es una plataforma de deliberaci\'{o}n c\'{i}vica con IA para escolares.
El aula se convierte en plaza digital: los estudiantes debaten problemas reales,
guiados por una IA que valida fuentes y modera el di\'{a}logo,
y producen propuestas ciudadanas concretas que llegan a quienes pueden actuar.
\textbf{\color{yteal}No ense\~{n}amos sobre ciudadan\'{i}a. La practicamos.}}
\end{tcolorbox}

\vspace{2pt}

%% ─── PROBLEMA / SOLUCIÓN ──────────────────────────────────────────────
\begin{tcolorbox}[
  enhanced,
  colback=ylightbg,
  colframe=yblue!25,
  arc=3mm,
  boxrule=0.8pt,
  left=8pt, right=8pt, top=5pt, bottom=5pt
]
\begin{minipage}[t]{0.465\textwidth}
  {\color{yorange}\fontsize{11}{14}\selectfont\textbf{El Problema}}\\[5pt]
  {\color{ygray}\footnotesize
    Los escolares peruanos carecen de espacios reales de participaci\'{o}n ciudadana.
    La educaci\'{o}n c\'{i}vica se reduce a memorizar conceptos, sin practicar el di\'{a}logo,
    la deliberaci\'{o}n ni el pensamiento cr\'{i}tico. La desinformaci\'{o}n digital
    los alcanza sin herramientas para evaluarla.}
\end{minipage}%
\hfill%
\textcolor{yblue!20}{\vrule width 0.6pt height 3cm}%
\hfill%
\begin{minipage}[t]{0.465\textwidth}
  {\color{yteal}\fontsize{11}{14}\selectfont\textbf{La Soluci\'{o}n}}\\[5pt]
  {\color{ygray}\footnotesize
    YariNET transforma el aula en un ecosistema participativo.
    Los estudiantes no estudian sobre ciudadan\'{i}a --- la ejercen.
    Debaten problemas reales de su barrio, con IA que guía el an\'{a}lisis
    de fuentes y la construcci\'{o}n de argumentos, generando propuestas
    que llegan a la comunidad.}
\end{minipage}
\end{tcolorbox}

\vspace{2pt}

%% ─── CÓMO FUNCIONA ────────────────────────────────────────────────────
\begin{tcolorbox}[
  enhanced,
  colback=yblue,
  colframe=yblue,
  arc=3mm,
  boxrule=0pt,
  left=8pt, right=8pt, top=8pt, bottom=10pt
]
{\color{ygold}\fontsize{10}{12}\selectfont\textbf{C\'{O}MO FUNCIONA}}\\[8pt]
\begin{center}
\begin{tikzpicture}[
  box/.style={
    rectangle, rounded corners=4pt,
    fill=white!12!yblue,
    draw=white!25!yblue,
    line width=0.5pt,
    text=white,
    font=\fontsize{8}{10}\selectfont\bfseries,
    minimum width=2.85cm,
    minimum height=0.85cm,
    align=center,
    inner sep=5pt
  },
  arr/.style={-Stealth, color=ygold, line width=1.3pt}
]
\node[box] (n1) {Profesor\\lanza reto};
\node[box, right=0.3cm of n1] (n2) {Alumnos\\investigan};
\node[box, right=0.3cm of n2] (n3) {IA modera\\el debate};
\node[box, right=0.3cm of n3] (n4) {Propuesta\\colectiva};
\node[box, right=0.3cm of n4] (n5) {Impacto en\\comunidad};
\draw[arr] (n1) -- (n2);
\draw[arr] (n2) -- (n3);
\draw[arr] (n3) -- (n4);
\draw[arr] (n4) -- (n5);
\end{tikzpicture}
\end{center}
\end{tcolorbox}

\newpage

%% ─── MINI HEADER PÁGINA 2 ────────────────────────────────────────────
\begin{tcolorbox}[
  enhanced,
  colback=yblue,
  colframe=yblue,
  arc=3mm,
  boxrule=0pt,
  left=10pt, right=10pt, top=5pt, bottom=5pt
]
\begin{minipage}[c]{0.08\textwidth}
  \includegraphics[width=\linewidth]{C:/Trabajo/CUMBRE_PLATFORM/Logotipos/YariNet.png}
\end{minipage}%
\hspace{10pt}%
\begin{minipage}[c]{0.88\textwidth}
  {\color{white}\fontsize{14}{18}\selectfont\textbf{YariNET}} \quad
  {\color{ygold}\fontsize{10}{12}\selectfont\textit{Piensa. Conecta. Act\'{u}a.}} \quad
  {\color{white!50!yblue}\fontsize{8}{10}\selectfont Educat\'{o}n Challenge 2026 · PUCP}
\end{minipage}
\end{tcolorbox}

\vspace{4pt}

%% ─── EJES EDUCATÓN 2026 ───────────────────────────────────────────────
\begin{tcolorbox}[
  enhanced,
  colback=ygold!10,
  colframe=ygold,
  arc=3mm,
  boxrule=1pt,
  left=8pt, right=8pt, top=5pt, bottom=5pt
]
{\color{yblue}\fontsize{10}{12}\selectfont\textbf{EJES EDUCAT\'{O}N CHALLENGE 2026}}\\[6pt]
\begin{minipage}[t]{0.465\textwidth}
  {\color{yblue}\small\textbf{Ciudadan\'{i}a Digital Activa}}\\[3pt]
  {\color{ygray}\footnotesize
    La IA gu\'{i}a a los estudiantes a buscar, analizar y validar
    informaci\'{o}n en tiempo real durante el debate --- desarrollando
    competencias cr\'{i}ticas aplicadas a decisiones ciudadanas reales.}
\end{minipage}%
\hfill%
\textcolor{ygold!50}{\vrule width 0.6pt height 2.5cm}%
\hfill%
\begin{minipage}[t]{0.465\textwidth}
  {\color{yblue}\small\textbf{Innovaci\'{o}n para la Participaci\'{o}n Ciudadana}}\\[3pt]
  {\color{ygray}\footnotesize
    YariNET usa IA para transformar el aula en un espacio real
    de deliberaci\'{o}n --- los estudiantes se convierten en agentes
    de cambio, con propuestas que conectan con autoridades y comunidad.}
\end{minipage}
\end{tcolorbox}

\vspace{2pt}

%% ─── TECNOLOGÍA + ROADMAP ─────────────────────────────────────────────
\begin{tcolorbox}[
  enhanced,
  colback=ylightbg,
  colframe=ypurple!35,
  arc=3mm,
  boxrule=0.8pt,
  left=8pt, right=8pt, top=5pt, bottom=5pt
]
\begin{minipage}[t]{0.465\textwidth}
  {\color{ypurple}\fontsize{10}{12}\selectfont\textbf{Tecnolog\'{i}a}}\\[5pt]
  {\color{ygray}\footnotesize
  \begin{itemize}[leftmargin=10pt, itemsep=0pt, topsep=0pt, parsep=0pt]
    \item IA conversacional en espa\~{n}ol
    \item Motor de validaci\'{o}n de fuentes en tiempo real
    \item Knowledge graph de conceptos c\'{i}vicos
    \item Dashboard docente con analytics de aula
    \item Accesible desde cualquier dispositivo
  \end{itemize}}
\end{minipage}%
\hfill%
\textcolor{ypurple!20}{\vrule width 0.6pt height 2.8cm}%
\hfill%
\begin{minipage}[t]{0.465\textwidth}
  {\color{ypurple}\fontsize{10}{12}\selectfont\textbf{Roadmap}}\\[5pt]
  {\color{ygray}\footnotesize
  \begin{itemize}[leftmargin=10pt, itemsep=0pt, topsep=0pt, parsep=0pt]
    \item \textbf{Fase 1:} Colegios Barrio PUCP (2026-2)
    \item \textbf{Fase 2:} Alianza MINEDU --- certificaci\'{o}n oficial
    \item \textbf{Fase 3:} Expansi\'{o}n nacional
    \item \textbf{Visi\'{o}n:} Edge AI para zonas sin internet (Jetson Nano)
  \end{itemize}}
\end{minipage}
\end{tcolorbox}

\vspace{2pt}

%% ─── FOOTER ───────────────────────────────────────────────────────────
\begin{center}
{\color{ygray}\fontsize{8.5}{10}\selectfont
\textit{Aprendemos\ $\cdot$\ Dialogamos\ $\cdot$\ Construimos juntos}
\quad $\bullet$ \quad
{\color{yblue}\textbf{Educat\'{o}n Challenge 2026}}
\quad $\bullet$ \quad
PUCP\ $\cdot$\ Barrio PUCP}
\end{center}

\end{document}
